%! Carrying Out a Test for the Difference of Two Population Means — Unit 7, Lesson 7.9 — Homework (Answer Key)
\documentclass[10pt]{article}
\usepackage{apstats-article}
\usepackage{apstats-key}   % pulls in -boxes; do NOT also load -boxes
\usepackage{grading}
\gradingassignment{apstats-u07-l09-hw}
\gradingcourse{advanced_statistics}
\newcommand{\TallMath}[1]{$\displaystyle #1\rule[-1.4em]{0pt}{3.2em}$}

\begin{document}

\pageheader{Unit 7, Lesson 7.9}{Homework --- Answer Key}
\namedateperiod

\vspace{0.1in}

\textit{Complete all four SPDC steps for each problem. Show all computations.}

\begin{enumerate}

\item \textbf{Study Hours.}

\begin{enumerate}[label=\alph*.]
  \item \textbf{State:} $H_0$: \ans{$\mu_1 - \mu_2 = 0$} \quad $H_a$: \ans{$\mu_1 - \mu_2 > 0$}\gradedkey{$H_0$:~\gradedblank[id=hw1a1, type=short, points=1, answer={{mu1 - mu2 = 0}}, check=text, display={\boldmath$\mu_1 - \mu_2 = 0$}]{3.8cm}{0.19in}\quad $H_a$:~\gradedblank[id=hw1a2, type=short, points=1, answer={{mu1 - mu2 > 0}}, check=text, display={\boldmath$\mu_1 - \mu_2 > 0$}]{3.8cm}{0.19in}}
  \item \textbf{Plan:} Method: \ans{Two-sample $t$-test for a difference of means}, conditions met: \ansline{Random samples; both $n \ge 30$ so CLT applies.}\gradedkey{\gradedblank[id=hw1b1, type=short, points=1, answer={{two-sample t-test for mu1 - mu2}}, check=text, display={\boldmath two-sample $t$-test for $\mu_1 - \mu_2$}]{4.8cm}{0.19in}\quad \gradedblank[id=hw1b2, type=rubric, points=1, display={\bfseries\boldmath random; CLT}]{1.25in}{0.19in}}
  \item \textbf{Do:} $SE = \sqrt{\dfrac{s_1^2}{n_1} + \dfrac{s_2^2}{n_2}} \approx$ \ans{1.39}
        \quad $t \approx$ \ans{1.80} \quad Technology: $p \approx$ \ans{0.037}, $df \approx$ \ans{59}\gradedkey{$SE=$~\gradedblank[id=hw1c1, type=short, points=1, answer=1.419, check=numeric, tolerance=0.01, display={\boldmath$1.419$}]{2.3cm}{0.19in}\quad $t=$~\gradedblank[id=hw1c2, type=short, points=1, answer=1.76, check=numeric, tolerance=0.02, display={\boldmath$1.76$}]{1.8cm}{0.19in}\quad $p=$~\gradedblank[id=hw1c3, type=short, points=1, answer=0.042, check=numeric, tolerance=0.003, display={\boldmath$0.042$}]{1.8cm}{0.19in}\quad $df=$~\gradedblank[id=hw1c4, type=short, points=1, answer=57.4, check=numeric, tolerance=0.6, display={\boldmath$57.4$}]{1.8cm}{0.19in}}
  \item \textbf{Conclude:} At $\alpha = 0.05$: \ansline{Because $p \approx 0.037 < \alpha = 0.05$, we reject $H_0$. We have convincing evidence that residential students study more hours per week on average than commuter students.}\gradedkeyline{\gradedblank[id=hw1d, type=rubric, points=2, display={\bfseries\boldmath $p \approx 0.042 < 0.05$: reject $H_0$; convincing evidence residential students study more on average.}]{\linewidth+6pt}{\gradedlinesheight{1}}}
\end{enumerate}

\vspace{0.05in}

\item \textbf{Reading Scores.}

\begin{enumerate}[label=\alph*.]
  \item \textbf{State:} $H_0$: \ans{$\mu_1 - \mu_2 = 0$} \quad $H_a$: \ans{$\mu_1 - \mu_2 \ne 0$}\gradedkey{$H_0$:~\gradedblank[id=hw2a1, type=short, points=1, answer={{mu1 - mu2 = 0}}, check=text, display={\boldmath$\mu_1 - \mu_2 = 0$}]{3.8cm}{0.19in}\quad $H_a$:~\gradedblank[id=hw2a2, type=short, points=1, answer={{mu1 - mu2 != 0}}, check=text, display={\boldmath$\mu_1 - \mu_2 \ne 0$}]{3.8cm}{0.19in}}
  \item \textbf{Plan:} Method: \ans{Two-sample $t$-test}. Conditions: \ansline{Random assignment; $n_1 = 40$ and $n_2 = 38 \ge 30$, CLT applies.}\gradedkey{\gradedblank[id=hw2b1, type=short, points=1, answer={{two-sample t-test for mu1 - mu2}}, check=text, display={\boldmath two-sample $t$-test for $\mu_1 - \mu_2$}]{4.8cm}{0.19in}\quad \gradedblank[id=hw2b2, type=rubric, points=1, display={\bfseries\boldmath random; $n$'s $\ge 30$}]{1.5in}{0.19in}}
  \item \textbf{Do:} $SE \approx$ \ans{2.31} \quad $t \approx$ \ans{1.56}
        \quad $p$ (two-sided) $\approx$ \ans{0.122} \quad $df \approx$ \ans{74}\gradedkey{$SE=$~\gradedblank[id=hw2c1, type=short, points=1, answer=2.388, check=numeric, tolerance=0.01, display={\boldmath$2.388$}]{2.3cm}{0.19in}\quad $t=$~\gradedblank[id=hw2c2, type=short, points=1, answer=1.51, check=numeric, tolerance=0.02, display={\boldmath$1.51$}]{1.8cm}{0.19in}\quad $p=$~\gradedblank[id=hw2c3, type=short, points=1, answer=0.136, check=numeric, tolerance=0.003, display={\boldmath$0.136$}]{1.8cm}{0.19in}\quad $df=$~\gradedblank[id=hw2c4, type=short, points=1, answer=73.5, check=numeric, tolerance=0.6, display={\boldmath$73.5$}]{1.8cm}{0.19in}}
  \item \textbf{Conclude:} \ansline{Because $p \approx 0.122 > \alpha = 0.05$, we fail to reject $H_0$. We do not have convincing evidence that the two literacy programs produce different mean reading scores.}\gradedkeyline{\gradedblank[id=hw2d, type=rubric, points=2, display={\bfseries\boldmath $p \approx 0.136 > 0.05$: fail to reject $H_0$; no convincing evidence the programs' mean scores differ.}]{\linewidth+6pt}{\gradedlinesheight{1}}}
\end{enumerate}

\vspace{0.05in}

\item \textbf{Recovery Time.}

\begin{enumerate}[label=\alph*.]
  \item \textbf{State:} $H_0$: \ans{$\mu_1 - \mu_2 = 0$} \quad $H_a$: \ans{$\mu_1 - \mu_2 < 0$}\gradedkey{$H_0$:~\gradedblank[id=hw3a1, type=short, points=1, answer={{mu1 - mu2 = 0}}, check=text, display={\boldmath$\mu_1 - \mu_2 = 0$}]{3.8cm}{0.19in}\quad $H_a$:~\gradedblank[id=hw3a2, type=short, points=1, answer={{mu1 - mu2 < 0}}, check=text, display={\boldmath$\mu_1 - \mu_2 < 0$}]{3.8cm}{0.19in}}

        \begin{teachernote}
        Group 1 = new protocol; Group 2 = standard. Since lower recovery time is better,
        $H_a: \mu_1 < \mu_2$, i.e., $\mu_1 - \mu_2 < 0$.
        \end{teachernote}

  \item \textbf{Plan:} Method: \ans{Two-sample $t$-test}. Conditions:\gradedkey{\gradedblank[id=hw3b1, type=short, points=1, answer={{two-sample t-test for mu1 - mu2}}, check=text, display={\boldmath two-sample $t$-test for $\mu_1 - \mu_2$}]{4.8cm}{0.19in}}
        \begin{itemize}
          \item Random: \ansline{Randomized experiment --- patients randomly assigned to protocols.}\gradedkey{\gradedblank[id=hw3b2, type=rubric, points=1, display={\bfseries\boldmath randomized experiment: random assignment --- met}]{4.6in}{0.19in}}
          \item Normal ($n < 30$): \ansline{Both distributions roughly symmetric, no outliers; normality condition met.}\gradedkey{\gradedblank[id=hw3b3, type=rubric, points=1, display={\bfseries\boldmath both roughly symmetric, no outliers --- met}]{4.6in}{0.19in}}
        \end{itemize}
  \item \textbf{Do:} $SE \approx$ \ans{1.21} \quad $t \approx$ \ans{$-2.73$}
        \quad Technology: $p \approx$ \ans{0.005} \quad $df \approx$ \ans{35}\gradedkey{$SE=$~\gradedblank[id=hw3c1, type=short, points=1, answer=1.249, check=numeric, tolerance=0.01, display={\boldmath$1.249$}]{2.3cm}{0.19in}\quad $t=$~\gradedblank[id=hw3c2, type=short, points=1, answer=-2.64, check=numeric, tolerance=0.02, display={\boldmath$-2.64$}]{1.8cm}{0.19in}\quad $p=$~\gradedblank[id=hw3c3, type=short, points=1, answer=0.006, check=numeric, tolerance=0.002, display={\boldmath$0.006$}]{1.8cm}{0.19in}\quad $df=$~\gradedblank[id=hw3c4, type=short, points=1, answer=36.0, check=numeric, tolerance=0.6, display={\boldmath$36.0$}]{1.8cm}{0.19in}}
  \item \textbf{Conclude:} At $\alpha = 0.05$: \ansline{Because $p \approx 0.005 < \alpha = 0.05$, we reject $H_0$. We have convincing evidence that the new stretching protocol reduces mean recovery time compared to the standard protocol.}\gradedkeyline{\gradedblank[id=hw3d, type=rubric, points=2, display={\bfseries\boldmath $p \approx 0.006 < 0.05$: reject $H_0$; convincing evidence the new protocol lowers mean recovery time.}]{\linewidth+6pt}{\gradedlinesheight{1}}}
\end{enumerate}

\vspace{0.05in}

\item \textbf{Fuel Efficiency.}

\begin{enumerate}[label=\alph*.]
  \item \textbf{State:} $H_0$: \ans{$\mu_1 - \mu_2 = 0$} \quad $H_a$: \ans{$\mu_1 - \mu_2 \ne 0$}\gradedkey{$H_0$:~\gradedblank[id=hw4a1, type=short, points=1, answer={{mu1 - mu2 = 0}}, check=text, display={\boldmath$\mu_1 - \mu_2 = 0$}]{3.8cm}{0.19in}\quad $H_a$:~\gradedblank[id=hw4a2, type=short, points=1, answer={{mu1 - mu2 != 0}}, check=text, display={\boldmath$\mu_1 - \mu_2 \ne 0$}]{3.8cm}{0.19in}}
  \item \textbf{Plan:} Method: \ans{Two-sample $t$-test}. Conditions: \ansline{Random samples; both dotplots roughly symmetric with no outliers.}\gradedkey{\gradedblank[id=hw4b1, type=short, points=1, answer={{two-sample t-test for mu1 - mu2}}, check=text, display={\boldmath two-sample $t$-test for $\mu_1 - \mu_2$}]{4.8cm}{0.19in}\quad \gradedblank[id=hw4b2, type=rubric, points=1, display={\bfseries\boldmath random; symmetric}]{1.5in}{0.19in}}
  \item \textbf{Do:} $SE \approx$ \ans{0.862} \quad $t \approx$ \ans{$-1.97$}
        \quad $p$ (two-sided) $\approx$ \ans{0.054}\gradedkey{$SE=$~\gradedblank[id=hw4c1, type=short, points=1, answer=0.850, check=numeric, tolerance=0.01, display={\boldmath$0.850$}]{2.3cm}{0.19in}\quad $t=$~\gradedblank[id=hw4c2, type=short, points=1, answer=-2.00, check=numeric, tolerance=0.02, display={\boldmath$-2.00$}]{1.8cm}{0.19in}\quad $p=$~\gradedblank[id=hw4c3, type=short, points=1, answer=0.051, check=numeric, tolerance=0.002, display={\boldmath$0.051$}]{1.8cm}{0.19in}}
  \item \textbf{Conclude:} \ansline{Because $p \approx 0.054 > \alpha = 0.05$, we fail to reject $H_0$. We do not have convincing evidence (at the 5\% level) that the two engine designs produce different mean fuel efficiencies.}\gradedkeyline{\gradedblank[id=hw4d, type=rubric, points=2, display={\bfseries\boldmath $p \approx 0.051 > 0.05$: fail to reject $H_0$; no convincing evidence the mean mpg differs.}]{\linewidth+6pt}{\gradedlinesheight{1}}}
\end{enumerate}

\end{enumerate}

\begin{reflectionbox}
A student says: ``We got $p = 0.18$, which means the two groups probably have equal means.''
Explain what is wrong with this statement. What does ``fail to reject $H_0$'' actually mean
when comparing two groups?
\ansline{The $p$-value does not measure the probability that $H_0$ is true. It is computed assuming $H_0$ is true. A large $p$-value means the data are consistent with the null, not that the null is correct.}
\ansline{``Fail to reject $H_0$'' means the data do not provide sufficient evidence against equal population means --- not that the means are proven equal.}
\ansline{}
\end{reflectionbox}

\begin{extensionbox}
\textbf{Preview of Lesson 7.10 --- Selecting an Inference Procedure.}
For each scenario below, identify the correct inference procedure (one-sample $t$-test,
matched-pairs $t$-test, or two-sample $t$-test) and explain your reasoning.
\begin{enumerate}[label=(\alph*)]
  \item Researchers measure blood pressure for the same 15 patients before and after
        taking a new medication.
  \item Two independent groups of students take different versions of a standardized test.
        Researchers compare mean scores.
  \item A factory samples 40 bolts and tests whether the mean diameter differs from the
        specified 10 mm.
\end{enumerate}
\ansline{(a) Matched-pairs $t$-test: measurements are paired (same subject, before and after).}
\ansline{(b) Two-sample $t$-test: two independent groups are compared.}
\ansline{(c) One-sample $t$-test: one sample compared to a known null value (10 mm).}
\ansline{}
\end{extensionbox}

\end{document}
